\scanpage{36}
\mspage{17}
et ceci posé, la relation (56) équivaut à
\[
(60)\qquad\left\{\begin{aligned}
\alpha'&=\zeta'\rho^{\deleted{\unsure{-1}}}(\zeta')^{-1}\\
\text{i.e.}\quad\eta\sigma_\infty(\eta)^{-1}&=(\zeta'\rho^{\deleted{\unsure{-1}}}(\zeta')^{-1})\sigma_\infty(\lambda)
\end{aligned}\right.
\]
\deleted{i.e. à}\quad\deleted{(61)}
En effet, la formule (75bis) exprimant $\alpha$ en termes de $\widetilde\alpha$
\[
(x)\qquad\alpha=\rho(\eta\widetilde\alpha\rho^{-1}(\eta)^{-1})
\]
se récrit, en passant aux inverses et posant
\[
(xx)\qquad\widetilde\alpha'=\widetilde\alpha^{-1}
\]
sous la forme
\[
(61)\qquad\uline{\alpha'}=\rho(\rho^{-1}(\eta)\widetilde\alpha'\eta^{-1})
=\uline{\eta\rho(\widetilde\alpha')\rho(\eta)^{-1}}
\]
\deleted{et la formule}
(On considère ici $\alpha,\alpha',\widetilde\alpha,\widetilde\alpha'$ comme déduits de $\lambda,\eta\in\widehat\Pi_{0,3}$ par la première relation (55) pour ce qui est de $\widetilde\alpha$, dont les autres se déduisent par $(x),(xx)$. D'ailleurs, la \inserted{2\textsuperscript{e}} formule (55) se\deleted{\unsure{t}} récrit alors, en passant aux inverses
\[
\widetilde\alpha'=\rho^{-1}(\zeta)\zeta^{-1}
\]
i.e.
\[
\rho(\widetilde\alpha')=\zeta\rho(\zeta)^{-1}
\]
\deleted{de sorte que (61), \unsure{(62)} s'écrit encore}
\inserted{qui est équivalent} \inserted{\unsure{(61)}}
\[
\alpha'=\eta\zeta\rho(\zeta^{-1}\eta^{-1})
\]
\deleted{i.e. $\alpha'$ à (60)},
i.e. à (60), avec
\[
\zeta'=\eta\zeta.
\]

\uline{En résumé} Les autom. du groupoïde $\overline{\underline{\mathrm{VI}}}_J$, induisant l'identité sur les sommets et \uline{commutant à l'action de $G_J$}, correspondent aux
\[
(62)\qquad\boxed{(\lambda,\eta,\zeta')\in\widehat\Pi_{0,3}\times\widehat\Pi_{0,3}\times\widehat\Pi_{0,3}}
\]
satisfaisant la relation des lacets ?
\margin{35}
